Clinical practice guidelines provide recommendations for the diagnosis, treatment, and monitoring of patients, but different organisations often present similar information in different ways. They may use different terminology, document structures, recommendation formats, and measurement units. In diabetes care, this can make it difficult to combine information from multiple guidelines and use it consistently for tasks such as clinical information retrieval and educational question generation.

This thesis investigates how diabetes guidelines from different organisations can be represented in a common knowledge graph. The aim is to bring clinical entities, recommendations, numerical targets, and their relationships into a structured form that can be queried across guideline sources. The work also examines whether the resulting graph can be used as a controlled basis for generating study questions whose answer-bearing content is grounded in the extracted information.

The proposed system consists of four stages: document ingestion, knowledge graph construction, normalisation, and graph-anchored question generation. It was evaluated using eight guideline documents from four organisations covering Type 1 and Type 2 diabetes. The resulting graph contains 2,503 nodes and 3,485 edges across 18 relation types. Normalisation of entities and measurement units produced eight cross-guideline equivalence relations between selected numerical targets. The GPT-4o-mini question-generation run produced 1,289 questions, of which 1,280 passed the implemented validation checks. The accepted catalogue contains 3,048 stored reformulations and 294 multiple-choice questions. A separate sample of 30 questions, taken from an earlier partial run with
Claude Opus 4.8, was reviewed by two medical doctors; three clear errors were traced to incorrect graph relations, while one additional question raised a difference in clinical practice rather than a factual error.

The results show that the knowledge graph can bring together information from differently structured diabetes guidelines and support queries across guideline sources. The graph can also provide a controlled basis for generating educational questions whose factual content is tied to specific graph relations. However, the results also show that the quality of the generated questions depends on the quality of the extracted graph, particularly for clinical relations. The framework is therefore intended as a foundation for applications such as cross-guideline comparison, structured clinical information retrieval, and educational question generation, rather than as a replacement for the original clinical guidelines.
